% ===== 原文 pp.106–124（PDF pp.113–131） =====
% [原文 p.106]
作为附注，讨论 $d_{31}, d'_{31}$ 的任意性。$d_{31}$ 的任意性由满足
\[
  d_{X_1}^{i+1}\circ f^i + f^{i+1}\circ d_{X_3}^i = 0
\]
的 $f$ 表示为 $d_{31}+f$。也就是说，$f$ 定义了复形的态射 $f\colon X_3\to X_1[1]$。这一部分带来任意性；当 $f$ 在 $K(\cC)$ 中确定不同的态射时，对应的 $X_{123}$ 一般不在 $K(\cC)$ 中同构。

如此，已知 $K(\cC)$ 构成三角范畴，而它的正合三角序列的结构以映射锥为基础。阿贝尔范畴 $\cC$ 中复形的映射锥所确定的 $\Comp(\cC)$ 的短正合列，可以看作 $K(\cC)$ 中的短正合列。另一方面，由复形的短正合列可得长正合列（引理 4.52）。用对应的函子 $K(\cC)\to\Ab$ 的语言来说，便得到以下内容。

\begin{proposition}{5.25}
设 $\cC$ 为阿贝尔范畴。函子 $H^0\colon K(\cC)\to\Ab$ 由对应
\[
  K(\cC)\ni X \mapsto H^0(X)\in\Ab
\]
确定。$H^0$ 构成上同调函子。
\end{proposition}

% [Tutor Note: 原文 p.106 命题 5.25 的「コホモロジー的関手」译作「上同调函子」（cohomological functor），指把正合三角序列映为长正合列的函子；由 $H^0$ 与平移函子 $\T$ 合成即得各 $H^i$。可验证锚点：对正合三角 $X\xrightarrow{u}Y\to Z\to \T X$ 应用 $H^0$，得长正合列 $\cdots\to H^0(X)\to H^0(Y)\to H^0(Z)\to H^0(\T X)=H^1(X)\to\cdots$，这正是正文所说的「由短正合列得长正合列（引理 4.52）」在函子语言下的表述。]

\section{三角范畴的局部化与导出范畴}

\subsection{三角范畴的局部化}

\begin{definition}{5.26}
设 $\cT$ 为三角范畴。$\Ob(\cT)$ 的子类 $\cN$ 在满足以下条件时称为零系（null system）。
\begin{itemize}
  \item $0\in\cN$。
  \item $X\in\cN$ 当且仅当 $\T(X)\in\cN$。
  \item 在 $\cT$ 的正合三角序列 $(X,Y,Z;u,v,w)$ 中，若 $X,Y\in\cN$，则 $Z\in\cN$。
\end{itemize}
\end{definition}

% [Tutor Note: 「ゼロ系」英文为 null system，中文暂译「零系」，即 Verdier 意义下三角范畴局部化的核。第三条配合 (Tr2) 的旋转即可推出「三者中任意两个属于 $\cN$ 则第三个也属于 $\cN$」：由第二条知 $\cN$ 对平移 $\T$ 封闭，对旋转后的三角序列 $(Y,Z,\T X)$ 施用第三条，当 $Y,Z\in\cN$ 时得 $\T X\in\cN$，再由第二条得 $X\in\cN$。（此性质正文未明言，后文推 $S(\cN)$ 的积闭性时要用。）]

\begin{definition}{5.27}
对三角范畴 $\cT$ 的零系 $\cN$，关于 $\cT$ 的态射定义
\[
  S(\cN) := \{\, f \mid X\xrightarrow{f}Y\to Z\to \T X,\ Z\in\cN \,\}.
\]
\end{definition}

\begin{proposition}{5.28}
$S(\cN)$ 构成范畴 $\cT$ 的积闭系（定义 4.54）。
\end{proposition}

\begin{proof}
利用 $\cT$ 中正合三角结构的公理 (Tr0)–(Tr5)，确认 $S(\cN)$ 满足 $\cT$ 的积闭系（定义 4.54）的条件 (S1)–(S4) 即可。(S1) 由 (Tr1) 得到，(S2) 由八面体公理 (Tr5) 得到，(S3) 由 (Tr4) 得到，(S4) 则通过两次使用 (Tr4) 而得到 [34, p43]。（作为习题正合适。）
\end{proof}

% [Tutor Note: 原文此处以「恰好是一道习题」省略了逐条验证，补出可核对的路标：(S1) 取正合三角 $X\xrightarrow{\id_X}X\to0\to\T X$，因 $0\in\cN$ 故 $\id_X\in S(\cN)$；(S4) 在加法范畴中（本书 $\cT$ 由定义 5.4 是加法范畴）等价于「存在 $s\in S(\cN)$ 使 $s\circ(f-g)=0$ $\Leftrightarrow$ 存在 $t\in S(\cN)$ 使 $(f-g)\circ t=0$」（此即定义 4.54 的 (S4) 加注意 4.55），其中两方向各需对 $f-g$ 的映射锥用一次 (Tr4) 作同型替换。]

将 $S(\cN)$ 的局部化（作为所得的加法范畴）记作 $\cT/\cN$。由对应 $f\mapsto[(f,\id_X)]$，$f\colon X\to Y$ 所确定的加法函子 $\epsilon\colon\cT\to\cT/\cN$ 是三角函子，这一点将在下面看到。

% [原文 p.107]
\begin{proposition}{5.29}
考虑三角范畴 $\cT$ 关于零系 $\cN$ 的局部化 $\epsilon\colon\cT\to\cT/\cN$。$\cT/\cN$ 的三角序列中，与 $\cT$ 的正合三角序列在 $\epsilon$ 下的像同构者全体构成 $\cT/\cN$ 的正合三角序列，由此 $\cT/\cN$ 又成为三角范畴。
\end{proposition}

\begin{proof}
只要确认 $\cT$ 的正合三角序列在 $\epsilon$ 下的像满足正合三角序列的公理（定义 5.4）的 (Tr0)–(Tr5) 即可。(Tr0) 是显然的，(Tr1) 可作为像
\[
  \epsilon\bigl(X\xrightarrow{\id_X}X\xrightarrow{0}0\xrightarrow{0}\T X\bigr)
\]
得到。(Tr2) 是因为：对应于 $\cT/\cN$ 的态射 $X\xrightarrow{\tilde u}Y$，存在 $X\xleftarrow{s}X'\xrightarrow{u}Y$，$s\in S(\cN)$，$u\in\Hom_{\cT}(X,Y)$，于是考虑像
\[
  \epsilon\bigl(X'\xrightarrow{u}Y\xrightarrow{v}C(u)\xrightarrow{w}\T X'\bigr),
\]
在 $\cT/\cN$ 中 $X$ 与 $X'$ 同构，故得所欲证。(Tr3) 到 (Tr5) 可由把 $\cT$ 中 (Tr3) 到 (Tr5) 的性质经函子 $\epsilon$ 搬运而得到。
$\epsilon$ 是三角函子这一点，由 $\cT/\cN$ 中正合三角序列的定义也显而易见。
\end{proof}

\begin{remark}{5.30}
由于 $\cT/\cN$ 中的正合三角序列已经确定，故由正合三角序列的公理 (Tr0)–(Tr5) 例如可知以下事项。首先，对 $s\in S(\cN)$，$\cT$ 的正合三角序列
\[
  X\xrightarrow{s}Y\to Z\to\T X
\]
（即 $Z\in\cN$）在 $S(\cN)^{-1}$ 下的像，由于 $s$ 被送到同构态射，故在 $\cT/\cN$ 中与
\[
  X\xrightarrow{\id_X}X\to0\to\T X
\]
同构。也就是说，$\cN$ 的对象在 $\cT/\cN$ 中成为零对象。
\end{remark}

由范畴的局部化所具有的泛性质（定理 4.58），可得以下结论。

\begin{corollary}{5.31}
对于三角范畴 $\cT$ 与零系 $\cN$，任意三角范畴 $\cT'$ 与三角函子 $F\colon\cT\to\cT'$，只要它把 $\cT$ 中 $S(\cN)$ 的态射送到 $\cT'$ 中的同构态射，则 $F$ 作为三角函子经过 $\epsilon\colon\cT\to\cT/\cN$ 分解。
\end{corollary}

% [Tutor Note: 注意 5.30 的「$\cN$ 的对象成为零对象」是下文 $D(\cC)=K(\cC)/\cN$ 里「无界复形可取零对象替换」的来源，具体验证可取 $s=\id_X$（此时三角序列 $X\xrightarrow{\id}X\to0\to\T X$ 已在 $\cT$ 中）与一般的 $s\in S(\cN)$ 对比；$\cT/\cN$ 中 $\epsilon(s)$ 为同构，故 $\epsilon(X)\cong\epsilon(X')$ 而 $Z\in\cN$ 被消去。]

\subsection{导出范畴（再）}

下面来看导出范畴 $D(\cC)$（定义 4.61）是三角范畴。

\begin{lemma}{5.32}
把阿贝尔范畴的同伦范畴 $K(\cC)$ 的零系定义为
\[
  \cN := \bigl\{\{(X^i,d^i)\}_{i\in\Z}\in K(\cC)\bigm| H^i(X)=0,\ \forall i\in\Z\bigr\}.
\]
则 $S(\cN)$ 与拟同构映射的全体一致。
\end{lemma}

\begin{proof}
把命题 5.25 的上同调函子 $H^0$ 应用于 $K(\cC)$ 的正合三角序列
\[
  X\xrightarrow{s}Y\to Z\to\T X,\qquad Z\in\cN
\]
便得到长正合列
% [原文 p.108]
\[\begin{tikzcd}[row sep=2.4em, column sep=3.4em]
\cdots & \cdots & \cdots \arrow[dll] \\
H^{i-1}(X) \arrow[r,"s^{i-1}"] & H^{i-1}(Y) \arrow[r] & 0 \arrow[dll] \\
H^{i}(X) \arrow[r,"s^{i}"] & H^{i}(Y) \arrow[r] & 0 \arrow[dll] \\
\cdots & \cdots & \cdots
\end{tikzcd}\]
由此，这里用到了对每个 $i\in\Z$ 有 $H^i(Z)=0$。于是对每个 $i$ 得到同构
\[
  H^i(X)\simeq H^i(Y).
\]
也就是说，证明了 $s$ 是拟同构映射。
\end{proof}

由此，引理 4.60 成立，也就是说，已知阿贝尔范畴 $\cC$ 的导出范畴 $D(\cC)$ 构成三角范畴。

以下命题给出 $D(\cC)$ 的正合三角序列的基本刻画。

\begin{proposition}{5.33}
设 $\cC$ 为阿贝尔范畴，考虑 $\Comp(\cC)$ 中的短正合列
\begin{equation}
  0\to X\xrightarrow{u}Y\xrightarrow{v}Z\to 0 \tag{5.14}
\end{equation}
此时，在 $D(\cC)$ 中存在正合三角序列
\[
  X\xrightarrow{u}Y\xrightarrow{v'}Z\xrightarrow{w'}\T(X).
\]
其中，$X, Y, Z$ 是把 (5.14) 中 $\Comp(\cC)$ 的对象原样视为 $D(\cC)$ 的对象，$u\colon X\to Y$ 也是由 (5.14) 中 $\Comp(\cC)$ 的态射 $u$ 诱导的 $D(\cC)$ 的态射。
\end{proposition}

\begin{proof}
首先，对 $\Comp(\cC)$ 中的态射 $u\colon X\to Y$，关于其映射锥 $C(u)$ 构造到 $Z$ 的拟同构映射 $s\colon C(u)\to Z$。令 $s:=\{s^n\}_{n\in\Z}$ 由
\begin{equation}
  s^n = \begin{pmatrix} v^n & 0 \end{pmatrix}\colon
  \begin{pmatrix} Y^n \\ (X[1])^n \end{pmatrix} \longrightarrow Z^n \tag{5.15}
\end{equation}
给出。它是复形的态射，这由 $d_Z\circ s-s\circ d_{C(u)}=0$ 可得。为了看出它是拟同构映射，考虑以下复形的短正合列。
\[
  0\to\Ker(s)\xrightarrow{i}C(u)\xrightarrow{s}Z\to 0.
\]
这里，$\Ker(s):=\{K^n,d_K^n\}$ 满足 $K^n:=\Ker(s^i\colon C(u)^i\to Z^i)$，并由此得到图表
\[\begin{tikzcd}[row sep=2.8em, column sep=3.0em]
K^n \arrow[r,"i^n"] \arrow[d,"d_K^n"] & C(u)^n \arrow[r] \arrow[d,"d_{C(u)}^n"] & Z^n \arrow[r] \arrow[d,"d_Z^n"] & 0 \\
K^{n+1} \arrow[r,"i^{n+1}"] & C(u)^{n+1} \arrow[r] & Z^{n+1} \arrow[r] & 0
\end{tikzcd}\]
% [Tutor Note: (5.15) 中 $(X[1])^n=X^{n+1}$，而 $C(u)^n=Y^n\oplus (X[1])^n$；因此 $s^n$ 就是「投影到 $Y^n$ 后取 $v^n$」的行向量。可核对 $d_Z\circ s-s\circ d_{C(u)}=0$：$C(u)$ 的微分矩阵为 $\begin{psmallmatrix} d_Y & u^{n+1}\\ 0 & d_{X[1]}\end{psmallmatrix}$（见 p.109 式 (5.16) 前后），与 $s^n=(v^n\ 0)$ 合成后得到 $v^{n+1}\circ d_Y^n-d_Z^n\circ v^n$ 与 $v^n\circ u^{n+1}$ 两项，前者为 $0$（$v$ 是复形态射），后者为 $0$（(5.14) 逐项正合，故 $v^{n+1}\circ u^{n+1}=0$）。]

% [原文 p.109]
其中，由 $d_{C(u)}^n\circ i^n$ 经过 $i^{n+1}\colon K^{n+1}\to C(u)^{n+1}$ 分解，把由此唯一确定的态射记作 $d_K^n\colon K^n\to K^{n+1}$。由图表的交换性可知，对每个 $n\in\Z$ 有 $d_K^{n+1}\circ d_K^n=0$，于是得到复形 $\Ker(s)$。

现在想证明对每个 $i\in\Z$ 有 $H^i(\Ker(s))=0$。若能做到这一点，则由命题 5.25 可知 $s$ 是拟同构映射。首先，注意 (5.15) 可得
\[
  K^i=\Ker(s^i\colon C(u)^i\to Z^i)=X^n\oplus (X[1])^n,
\]
\[
  i^n=\begin{pmatrix} u^n & 0 \\ 0 & \id_{(X[1])^n}\end{pmatrix}\colon
  \begin{pmatrix} X^n \\ (X[1])^n\end{pmatrix}\longrightarrow
  \begin{pmatrix} Y^n \\ (X[1])^n\end{pmatrix}
\]
可知。这里，若令
\begin{equation}
  d_K^n=\begin{pmatrix} d_X^n & \id_{(X[1])^n} \\ 0 & d_{X[1]}^n\end{pmatrix} \tag{5.16}
\end{equation}
则
\begin{align*}
  d_{C(u)}^n\circ i^n
  &= \begin{pmatrix} d_Y^n & u^{n+1} \\ 0 & d_{X[1]}^n\end{pmatrix}
     \begin{pmatrix} u^n & 0 \\ 0 & \id_{X^{n+1}}\end{pmatrix}
   = \begin{pmatrix} d_Y^n\circ u^n & u^{n+1} \\ 0 & d_{X[1]}^n\end{pmatrix}\\
  &= \begin{pmatrix} u^{n+1}\circ d_X^n & u^{n+1} \\ 0 & d_{X[1]}^n\end{pmatrix}
   = \begin{pmatrix} u^{n+1} & 0 \\ 0 & \id_{X^{n+2}}\end{pmatrix}
     \begin{pmatrix} d_X^n & \id_{X^{n+1}} \\ 0 & d_{X[1]}^n\end{pmatrix}
   = i^{n+1}\circ d_K^n
\end{align*}
可知 $\Ker(s)=C(\id_X)$。$C(\id_X)$ 在 $\Comp(\cC)$ 中与零对象同伦等价，这一点已在 5.2 节 [(Tr1)] 中看到；也就是说，$C(\id_X)$ 与零对象具有相同的上同调，故 $H^i(C(\id_X))=0$ 得证。
\end{proof}

% [Tutor Note: 底本 p.109 式 (5.16) 直后的计算中，第一个矩阵右上角印作 $u^n$；据作者正誤表（2014-05-20 版）p.109 line 11 订正为 $u^{n+1}$，此处按订正后的 $u^{n+1}$ 排印。核对依据：映射锥在次数 $n$ 处的微分为 $d_{C(u)}^n(y,x)=(d_Y^n y+u^{n+1}x,\ d_{X[1]}^n x)$，即矩阵 $\begin{psmallmatrix} d_Y^n & u^{n+1}\\ 0 & d_{X[1]}^n\end{psmallmatrix}$；而同一计算右端确实出现 $u^{n+1}$，故原印 $u^n$ 为笔误。另注：原文此句左侧印作 $K^i=\cdots=X^n\oplus (X[1])^n$，把 $i$ 与 $n$ 混用，译文照底本录入；上下文（以及紧接着的 $i^n$）表明这里应统一读作次数 $n$。]

\section{投射分解与导出范畴（再）}

对于阿贝尔范畴 $\cC$ 的上方有界导出范畴 $D^-(\cC)$，由上方有界的投射对象构成的同伦范畴 $K^-(\cP)$ 成为全子范畴。特别地，当 $\cC$ 具有足够多的投射对象时（定义 4.64），下面证明可得三角等价
\[
  D^-(\cC)\simeq K^-(\cP).
\]
特别地，关于（上下）有界的导出范畴、同伦范畴，$D^b(\cC)\simeq K^b(\cP)$ 成立（定理 5.46）。

% [Tutor Note: 底本 p.109 此句印作「$C$ が十分射影対象と持つとき」，据作者正誤表（2014-05-20 版）p.109 line −10 订正为「$C$ が十分射影対象を持つとき」（助词「と」→「を」）。中译「具有足够多的投射对象」在两种印法下语义相同。「十分射影対象を持つ」（has enough projective objects）与「各对象都有投射分解」等价，即本书定义 4.64 所采用的用法。]

\subsection{阿贝尔范畴中投射分解的基本事实}

在 4.8.1 节中，讨论了若阿贝尔范畴 $\cC$ 含有足够多的投射对象，则各对象 $X$ 具有投射分解。当 $X, Y\in\cC$ 具有投射分解时，关于态射 $f\colon X\to Y$，有以下结论（cf. [11], [17] 等）。

\begin{lemma}{5.34}
设阿贝尔范畴 $\cC$ 中的态射 $f\colon X\to Y$ 与 $X, Y\in\cC$ 的投射分解
% [原文 p.110]
\[
  0\leftarrow X\xleftarrow{p_0}P_0\xleftarrow{p_1}P_1\leftarrow\cdots,\qquad
  0\leftarrow Y\xleftarrow{q_0}Q_0\xleftarrow{q_1}Q_1\leftarrow\cdots
\]
给出。此时，
\begin{enumerate}[label=(\roman*)]
  \item 存在使以下图表交换的态射组 $\{\phi_i\colon P_i\to Q_i\}_{i=0,1,2,\ldots}$。
\[\begin{tikzcd}[row sep=2.4em, column sep=2.5em]
\cdots \arrow[r,"p_3"] & P_2 \arrow[r,"p_2"] \arrow[d,"\phi_2"] & P_1 \arrow[r,"p_1"] \arrow[d,"\phi_1"] & P_0 \arrow[r,"p_0"] \arrow[d,"\phi_0"] & X \arrow[r,"0"] \arrow[d,"f"] & 0 \arrow[d] \\
\cdots \arrow[r,"q_3"] & Q_2 \arrow[r,"q_2"] & Q_1 \arrow[r,"q_1"] & Q_0 \arrow[r,"q_0"] & Y \arrow[r,"0"] & 0
\end{tikzcd}\]
  \item 把这样的态射组 $\{\phi_i\}_{i=0,1,2,\ldots}$ 与 $f$ 合并视为复形的态射时，$\{\phi_i\}_{i=0,1,2,\ldots}$ 除同伦等价外唯一。
\end{enumerate}
\end{lemma}

% [Tutor Note: 底本 p.110 该图式上、下两行的左端印作「$\cdots\to0\to P_2\xrightarrow{p_2}\cdots$」「$\cdots\to0\to Q_2\xrightarrow{q_2}\cdots$」；据作者正誤表（2014-05-20 版）p.110「図式」条订正为「$\cdots\xrightarrow{p_3}P_2\xrightarrow{p_2}\cdots$」「$\cdots\xrightarrow{q_3}Q_2\xrightarrow{q_2}\cdots$」，此处按订正后排印。核对：由 $0\leftarrow X\xleftarrow{p_0}P_0\xleftarrow{p_1}P_1\leftarrow\cdots$ 知 $p_i\colon P_i\to P_{i-1}$，故射入 $P_2$ 的应是 $p_3\colon P_3\to P_2$，不能是零射。]

\begin{proof}
(i) 由于 $P_0$ 是投射的，且 $Q_0\to Y$ 是满射，态射 $f\circ p_0\colon P_0\to Y$ 可提升为 $\phi_0\colon P_0\to Q_0$。其次，由 $p_0\circ p_1=0$ 得
\[
  q_0\circ(\phi_0\circ p_1)=(q_0\circ\phi_0)\circ p_1=(f\circ p_0)\circ p_1=f\circ(p_0\circ p_1)=0.
\]
因此 $\phi_0\circ p_1$ 经过 $\Ker(q_0)$ 分解；但由 $Q_1\xrightarrow{q_1}Q_0\xrightarrow{q_0}Y$ 的正合性有 $\Ker(q_0)\simeq\Img(q_1)$，故 $\phi_0\circ p_1$ 经过 $\Img(q_1)$ 分解。也就是说，存在态射 $P_1\to\Img(q_1)$，它与 $\Img(q_1)\to Q_0$ 的合成等于 $\phi_0\circ p_1$。另一方面，$q_1\colon Q_1\to Q_0$ 也经过 $\Img(q_1)$ 分解，且该态射 $Q_1\to\Img(q_1)$ 是满射。由于 $P_1$ 是投射的，态射 $P_1\to\Img(q_1)$ 可提升为 $P_1\to Q_1$。把它记作 $\phi_1\colon P_1\to Q_1$。重复此过程便得到 $\phi_2,\phi_3,\ldots$。

(ii) 设 $\{\phi_i\}_{i=0,1,2,\ldots}$ 与 $\{\phi'_i\}_{i=0,1,2,\ldots}$ 都满足上述 (i) 的条件。此时，由态射的交换性
\[
  \bigl(P_0\xrightarrow{p_0}X\xrightarrow{f}Y\bigr)
  =\bigl(P_0\xrightarrow{\phi_0}Q_0\xrightarrow{q_0}Y\bigr)
\]
有 $q_0\circ(\phi'_0-\phi_0)=0$，故 $\phi'_0-\phi_0$ 经过 $\Ker(q_0)\simeq\Img(q_1)$ 分解。又因 $Q_1\to\Img(q_1)$ 是满射，由 $P_0$ 的投射性，态射 $P_0\to\Img(q_1)$ 可提升为 $h_0\colon P_0\to Q_1$，且 $\phi'_0-\phi_0=q_1\circ h_0$。
其次，由态射的交换性
\[
  \bigl(P_1\xrightarrow{p_1}P_0\xrightarrow{\phi_0}Q_0\bigr)
  =\bigl(P_1\xrightarrow{\phi_1}Q_1\xrightarrow{q_1}Q_0\bigr)
\]
有 $q_1\circ(\phi'_1-\phi_1)=(\phi'_0-\phi_0)\circ p_1=q_1\circ h_0\circ p_1$，即得 $q_1\circ\bigl((\phi'_1-\phi_1)-h_0\circ p_1\bigr)=0$。因此，与上面同样的论证给出存在 $h_1\colon P_1\to Q_2$ 使 $(\phi'_1-\phi_1)-h_0\circ p_1=q_2\circ h_1$；改写此等式得
\[
  \phi'_1-\phi_1=h_0\circ p_1+q_2\circ h_1.
\]
同样地重复此过程，对每个 $i$ 得到
% [原文 p.111]
\[
  \phi'_i-\phi_i=h_{i-1}\circ p_i+q_{i+1}\circ h_i.
\]
由此得到，也就是说，由 $f$ 与 $\{\phi_i\}$ 构成的复形的态射和由 $f$ 与 $\{\phi'_i\}$ 构成的复形的态射是同伦等价的。
\end{proof}

% [Tutor Note: 底本 p.110 证明 (i) 首行印作「$f\circ p_0\colon P\to Y$」；据作者正誤表（2014-05-20 版）p.110「証明 (i) line 1」订正为「$f\circ p_0\colon P_0\to Y$」，此处按订正后排印（因 $p_0\colon P_0\to X$）。]

如 (4.18) 中所讨论的，对 $\cC$ 中的态射 $f\colon X\to Y$，把 $X, Y$ 视为复形，并设对应于投射分解
\[
  0\leftarrow X\xleftarrow{p_0}P_0\xleftarrow{p_1}P_1\leftarrow\cdots,\qquad
  0\leftarrow Y\xleftarrow{q_0}Q_0\xleftarrow{q_1}Q_1\leftarrow\cdots
\]
的由投射对象构成的复形为 $P, Q$，把 $p_0, q_0$ 所确定的复形的态射记作 $p\colon P\to X$，$q\colon Q\to Y$。此时，引理 5.34 的断言等价于：存在复形的态射 $\phi\colon P\to Q$ 满足 $q\circ\phi=\phi\circ p$，且除同伦外唯一。

% [Tutor Note: 底本 p.111 此处印作「$q\circ\phi=\phi\circ p$」；按态射的类型 $p\colon P\to X$、$q\colon Q\to Y$、$\phi\colon P\to Q$，$\phi\circ p$ 无法合成，正确的交换关系应为 $q\circ\phi=f\circ p$（见下文定理 5.42 的图式与 $K^-(\cC)$ 中可换的要求）。作者正誤表未收此条，故译文照底本录入，并在此注明。可核对：$q\circ\phi,\ f\circ p$ 都是 $P\to Y$ 的复形态射。]

\begin{corollary}{5.35}
阿贝尔范畴 $\cC$ 的对象 $X\in\cC$ 的投射分解，若存在则除同伦等价外唯一确定。也就是说，对于两个投射分解
\[
  0\leftarrow X\xleftarrow{p_0}P_0\xleftarrow{p_1}P_1\leftarrow\cdots,\qquad
  0\leftarrow X\xleftarrow{q_0}Q_0\xleftarrow{q_1}Q_1\leftarrow\cdots,
\]
对应的投射对象的复形 $P, Q$ 是同伦等价的。
\end{corollary}

\begin{proof}
在引理 5.34 中取 $Y=X$、$f=\id_X$ 的情形。此时得到复形的态射 $\phi\colon P\to Q$，$\psi\colon Q\to P$。于是得到复形的态射 $\psi\circ\phi\colon P\to P$，$\phi\circ\psi\colon Q\to Q$，由引理 5.34(ii) 它们分别与 $\id_P, \id_Q$ 同伦。
\end{proof}

\subsection{投射分解的同调扰动理论}

其次，考虑把复形替换为由投射对象构成的复形。对于阿贝尔范畴 $\cC$，以 $\Comp(\cC)$ 表示 $\cC$ 的复形构成的（阿贝尔）范畴，以 $K(\cC)$ 表示其同伦范畴。另一方面，设 $\cP$ 为仅由投射对象构成的 $\cC$ 的全子范畴，并把 $\Comp(\cP)\subset\Comp(\cC)$、$K(\cP)\subset K(\cC)$ 作为对应的全子范畴。本章将证明以下定理。

\begin{theorem}{5.36}
设阿贝尔范畴 $\cC$ 具有足够多的投射对象。此时，对于任意上方有界复形 $X\in\Comp^-(\cC)$，存在由投射对象构成的上方有界复形 $P\in\Comp^-(\cP)$ 与拟同构映射 $P\to X$。
\end{theorem}

% [Tutor Note: 「ホモロジー的摂動理論」译作「同调扰动理论」（homological perturbation theory），英文首次出现于下页注意 5.41。底本 p.111 定理 5.36 印作「アーベル圏 $C$ が十分射影的である」，据作者正誤表（2014-05-20 版）p.111「定理5.36 (etc.)」条订正为「アーベル圏 $C$ が十分射影対象を持つ」（同类订正也适用于本节其他处，见 p.116、p.119 正文），此处按订正后排印。]

首先，对 $X:=\{X^i,d_X^i\}_{i\in\Z}$，为每个 $X^n$ 取一个投射分解 $0\leftarrow X^n\xleftarrow{p_0^n}P_0^n\xleftarrow{p_1^n}P_1^n\leftarrow\cdots$。由引理 5.34，对每个 $n$ 存在态射组 $\{\phi_i^n\colon P_i^n\to P_i^{n+1}\}_{i=0,1,\ldots}$。于是得到以下双系列（交换图表）。
\begin{equation}
\begin{tikzcd}[row sep=2.2em, column sep=2.6em]
\vdots & \vdots \arrow[d,"p_2^{n-1}"] & \vdots \arrow[d,"p_2^{n}"] & \vdots \arrow[d,"p_2^{n+1}"] & \vdots \\
\cdots \arrow[r,"\phi_1^{n-2}"] & P_1^{n-1} \arrow[r,"\phi_1^{n-1}"] \arrow[d,"p_1^{n-1}"] & P_1^{n} \arrow[r,"\phi_1^{n}"] \arrow[d,"p_1^{n}"] & P_1^{n+1} \arrow[r,"\phi_1^{n+1}"] \arrow[d,"p_1^{n+1}"] & \cdots \\
\cdots \arrow[r,"\phi_0^{n-2}"] & P_0^{n-1} \arrow[r,"\phi_0^{n-1}"] \arrow[d,"p_0^{n-1}"] & P_0^{n} \arrow[r,"\phi_0^{n}"] \arrow[d,"p_0^{n}"] & P_0^{n+1} \arrow[r,"\phi_0^{n+1}"] \arrow[d,"p_0^{n+1}"] & \cdots \\
\cdots \arrow[r,"d_X^{n-2}"] & X^{n-1} \arrow[r,"d_X^{n-1}"] & X^{n} \arrow[r,"d_X^{n}"] & X^{n+1} \arrow[r,"d_X^{n+1}"] & \cdots
\end{tikzcd} \tag{5.17}
\end{equation}

% [原文 p.112]
不过，一般地 $\phi_i^n\circ\phi_i^{n-1}=0$ 并不成立，这一点要注意。
令 $P^n:=\bigoplus_{i=0}^\infty P_i^{n+i}$，并把 $d_0^n,d_1^n\colon P^n\to P^{n+1}$ 记作
\[
  d_0^n:=\bigoplus_{i=1}^\infty p_i^{n+i},\qquad
  d_1^n:=(-1)^n\bigoplus_{i=0}^\infty\phi_i^{n+i}.
\]
以下频繁省略 $n$，把 $d_0^n,d_1^n$ 简记为 $d_0,d_1$。把上面的交换图表「扭歪」画成如下图形，目标会更清楚。
\begin{equation}
\begin{tikzcd}[row sep=2.9em, column sep=2.6em]
\cdots \arrow[r] & P_2^{n+1} \arrow[r,"p_2^{n+1}"] & P_1^{n+1} \arrow[r,"p_1^{n+1}"] & P_0^{n+1} \arrow[ddd,"p_0^{n+1}"] & \cdots \\
P_2^{n} \arrow[r,"p_2^{n}"] \arrow[ur,"\phi_2^{n}"] & P_1^{n} \arrow[r,"p_1^{n}"] \arrow[ur,"\phi_1^{n}"] & P_0^{n} \arrow[ur,"\phi_0^{n}"] \arrow[dd,"p_0^{n}"] & & \\
P_1^{n-1} \arrow[r,"p_1^{n-1}"] \arrow[ur,"\phi_1^{n-1}"] & P_0^{n-1} \arrow[ur,"\phi_0^{n-1}"] \arrow[d,"p_0^{n-1}"] & & & \\
\cdots \arrow[r,"d_X^{n-2}"] & X^{n-1} \arrow[r,"d_X^{n-1}"] & X^{n} \arrow[r,"d_X^{n}"] & X^{n+1} \arrow[r,"d_X^{n+1}"] & \cdots
\end{tikzcd} \tag{5.18}
\end{equation}

% [Tutor Note: 底本 p.112 图式 (5.18) 中，从 $P_1^n$ 指向 $P_1^{n+1}$ 的斜箭头旁印作 $\phi_1^{n-1}$；据作者正誤表（2014-05-20 版）p.112「(5.18)」条，该处应为 $\phi_1^{n}$，此处按订正后排印。（同一图式中从 $P_1^{n-1}$ 指向 $P_1^n$ 的斜箭头才标 $\phi_1^{n-1}$。）可核对：$\phi_i^n\colon P_i^n\to P_i^{n+1}$，故「$n$ 层指向 $n+1$ 层」的斜箭头一律标 $\phi_i^{n}$。]

对每个 $n\in\Z$，以 $\pi^n\colon P^n\to P_0^n$ 为自然投影（满射）。关于图表 (5.17)（或图表 5.18），证明以下命题 5.37, 5.38。

\begin{proposition}{5.37}
对每个 $n\in\Z$，存在态射 $d_P^n\colon P^n\to P^{n+1}$，满足
\begin{itemize}
  \item $d_P$ 的 $P_i^{n+i}\to P_{i-1}^{n+i}$ 的成分是 $d_0^n$，
  \item $d_P$ 的 $P_i^{n+i}\to P_i^{n+i+1}$ 的成分是 $d_1^n$，且
  \item $\{P^n,d_P^n\}_{n\in\Z}$ 构成复形。
\end{itemize}
\end{proposition}

\begin{proposition}{5.38}
$\{p_0^n\circ\pi^n\colon P^n\to X^n\}_{n\in\Z}$ 是复形的态射，特别地是拟同构映射。
\end{proposition}

\begin{remark}{5.39}
在这些证明中常常对对象 $X$ 取从投射对象出发的满射 $P\to X$ 来讨论，这也可以看作是用函子 $\Hom_{\cC}(P,\ )$ 的像来代替阿贝尔范畴 $\cC$ 进行讨论。
\end{remark}

\begin{proof}[命题 5.37 的证明]
$d_P^n=d_P\colon P^n\to P^{n+1}$ 的存在性用归纳法证明。态射 $d_P$ 的 $P_i^{n+i}\to P_{i-1+j}^{n+i+j}$ 的成分在 $j=0$ 时为 $d_0$、在 $j=1$ 时为 $d_1$；此外还要说明：对于 $j\ge 2$，可以构造 $d_j\colon P_i^{n+i}\to P_{i-1+j}^{n+i+j}$，使得 $d_P:=d_0+d_1+d_2+\cdots$ 满足 $(d_P)^2=0$。现在 $d_0$ 与 $d_1$ 满足
\[
  d_0^2=0,\qquad d_0 d_1+d_1 d_0=0,
\]
但一般地 $(d_1)^2\neq 0$，这一点要注意。因此，关于 $d_0+d_1\colon P^n\to P^{n+1}$，有
% [原文 p.113]
\[
  (d_0+d_1)^2=(d_1)^2
\]
成立。也就是说，$(d_0+d_1)^2\colon P^n\to P^{n+2}$ 的 $P_i^{n+i}\to P_{i-2+j}^{n+i+j}$ 的成分在 $j=0,1$ 时是平凡的，但在 $j=2$ 时为 $(d_1)^2$。同样地，设 $d_j\colon P_i^{n+i}\to P_{i-1+j}^{n+i+j}$ 已对 $j=0,1,\ldots,k$ 构造出，且
\[
  (\bar d_k)^2\colon P^n\to P^{n+2},\qquad \bar d_k:=d_0+d_1+\cdots+d_k
\]
的 $P_i^{n+i}\to P_{i-2+j}^{n+i+j}$ 的成分 $\rho_j$ 在 $j=0,1,\ldots,k$ 时是平凡的。也就是说，设
\[
  (\bar d_k)^2=\rho_{k+1}+\rho_{k+2}+\cdots.
\]
此时，
\begin{align*}
  0 &= \bar d_k\circ(\bar d_k)^2-(\bar d_k)^2\circ\bar d_k\\
    &= \bar d_k\circ(\rho_{k+1}+\rho_{k+2}+\cdots)-(\rho_{k+1}+\rho_{k+2}+\cdots)\circ\bar d_k
\end{align*}
的 $P_i^{n+i}\to P_{i-3+j}^{n+i+j}$、$j=k+1$ 的成分来关注，便得到
\[
  0=d_0\circ\rho_{k+1}-\rho_{k+1}\circ d_0.
\]
此时，要说明存在 $d_{k+1}\colon P_i^{n+i}\to P_{i-1+(k+1)}^{n+i+(k+1)}$ 使 $\rho_{k+1}=d_0\circ d_{k+1}+d_{k+1}\circ d_0$。

% [Tutor Note: 底本 p.112 末～p.113 的指标序列有一串正誤表订正（2014-05-20 版：p112 line −6,−7；p113 line 2,4,6,11,14；p114 line −9），一律把 $P_{i-1-j}^{n+i+j}$ 型改作 $P_{i-1+j}^{n+i+j}$ 型（即把「$-j$」订正为「$+j$」），译文已全部按订正后排印。核对：$d_j$ 的双次数为 $(\Delta n,\Delta i)=(j,j-1)$，故 $d_j\colon P_i^{n+i}\to P_{i+j-1}^{n+i+j}$；取 $j=0,1$ 即得 $d_0\colon P_i^{n+i}\to P_{i-1}^{n+i}$（即 $p$ 型）与 $d_1\colon P_i^{n+i}\to P_i^{n+i+1}$（即 $\phi$ 型），与命题 5.37 的两条「成分」条件一致。]

\begin{lemma}{5.40}
当态射 $\rho\colon P\to P$ 满足 $d_0\circ\rho\mp\rho\circ d_0=0$ 时，存在满足
\[
  \rho=d_0\circ h\pm h\circ d_0
\]
的态射 $h\colon P\to P$（符号按同序对应）。
\end{lemma}

\begin{proof}
把 $P=\prod_n P^n=\prod_{m,n}P_m^{n+m}$ 分解时，可以分解为 $\rho=\sum_{k,l\in\Z}\rho_{k,l}$。这里 $\rho_{k,l}$ 是 $\rho$ 中对应于 $\rho_{k,l}\colon P_m^{n+m}\to P_{m+l}^{(n+k)+(m+l)}$ 的成分。当 $d_0\circ\rho\mp\rho\circ d_0=0$ 时，关于 $k,l$ 各自独立地有 $d_0\circ\rho_{k,l}\mp\rho_{k,l}\circ d_0=0$，故对每个 $\rho_{k,l}$ 证明命题即可。
取 $l>0$，固定 $n\in\Z$ 来考虑。
\[
  0=(d_0\circ\rho_{k,l}\mp\rho_{k,l}\circ d_0)|_{P_0^n}=d_0\circ\rho_{k,l}|_{P_0^n}
\]
因此，由 $P_0^n$ 的投射性，$\rho_{k,l}|_{P_0^n}\colon P_0^n\to P_l^{n+k}$ 经过 $P_{l+1}^{n+k}$ 分解。把对应的态射记作 $h^0\colon P_0^n\to P_{l+1}^{n+k}$。也就是说，得到 $\rho_{k,l}|_{P_0^n}=d_0\circ h^0$。
\[\begin{tikzcd}[row sep=2.8em, column sep=3.2em]
 & P_{l+1}^{n+k} \arrow[d,"d_0"]\\
P_0^n \arrow[ur,"h^0"] \arrow[r,"\rho_{k,l}"'] & P_l^{n+k}
\end{tikzcd}\]
% [原文 p.114]
其次，由 $0=(d_0\circ\rho_{k,l}\mp\rho_{k,l}\circ d_0)|_{P_1^n}$ 得
\[
  d_0\circ(\rho_{k,l}\mp h^0\circ d_0)|_{P_1^n}=0,
\]
故 $(\rho_{k,l}\mp h^0\circ d_0)|_{P_1^n}$ 经过 $P_{l+2}^{n+k}$ 分解。把对应的态射记作 $h^1\colon P_1^n\to P_{l+2}^{n+k}$。于是 $(\rho_{k,l}\mp h^0\circ d_0)|_{P_1^n}=d_0\circ h^1$，即移项得
\[
  \rho_{k,l}|_{P_1^n}=(d_0\circ h^1\pm h^0\circ d_0)|_{P_1^n}.
\]
到此为止所得的内容可由以下图表总结。
\[\begin{tikzcd}[row sep=2.8em, column sep=3.6em]
 & P_{l+2}^{n+k} \arrow[d,"d_0"]\\
P_1^n \arrow[ur,"h^1"] \arrow[r,"\rho_{k,l}"] \arrow[d,"d_0"] & P_{l+1}^{n+k} \arrow[d,"d_0"]\\
P_0^n \arrow[ur,"h^0"] \arrow[r,"\rho_{k,l}"'] & P_l^{n+k}
\end{tikzcd}\]
将此重复，对 $i\ge 0$ 得到
\[
  \rho_{k,l}|_{P_{i+1}^n}=(d_0\circ h^{i+1}\pm h^i\circ d_0)|_{P_{i+1}^n}.
\]
令 $h:=\sum_{i=0}^\infty h^i$，便得到
\[
  \rho_{k,l}=d_0\circ h+h\circ d_0.
\]
\end{proof}

回到命题 5.37 的证明。现在由补題 5.40 可取 $d_{k+1}$ 使 $\rho_k=d_0\circ d_{k+1}+d_{k+1}\circ d_0$。令 $\bar d_{k+1}:=\bar d_k+d_{k+1}$，则
\[
  (\bar d_{k+1})^2=(\bar d_k)^2+(\bar d_k\circ d_{k+1}+d_{k+1}\circ\bar d_k)+(d_{k+1})^2
\]
成立，其 $P_i^{n+i}\to P_{i-2+j}^{n+i+j}$ 的成分在 $j=0,\ldots,k$ 时为平凡（含 $d_{k+1}$ 的项不起作用），在 $j=k+1$ 时也有 $\rho_{k+1}+(d_0\circ d_{k+1}+d_{k+1}\circ d_0)=0$。以上完成了归纳法，命题 5.37 的证明到此结束。
\end{proof}

% [Tutor Note: 原文此处紧接「[命題5.38]の証明．」，该证明一直延续到 p120 才以 □ 收尾，跨本块多次落盘。为不使 chunk 末尾留下未闭合的 proof 环境，这一条证明按原文体例用 \para 手工起头（与 proof 环境同为「粗体・证明。」起排）；底本在 p117「$H^n(p)$ 的全射性」论证结束时另有一个 □，译文照排；p120 的单射性结束时再补最后一个 □。]

\para{命题 5.38 的证明。}
$\{p^n:=p_0^n\circ\pi^n\colon P^n\to X^n\}_{n\in\Z}$ 构成复形的态射，这一点容易看出。实际上，对每个 $n\in\Z$ 有
\begin{align*}
  (p_0^{n+1}\circ\pi^{n+1})\circ d_P^n
  &= p_0^{n+1}\circ(\pi^{n+1}\circ(d_0+d_1))\\
  &= p_0^{n+1}\circ(\pi^{n+1}\circ d_1)\\
  &= p_0^{n+1}\circ(d_1\circ\pi^n) = (p_0^{n+1}\circ d_1)\circ\pi^n\\
  &= (d_X^n\circ p_0^n)\circ\pi^n = d_X^n\circ(p_0^n\circ\pi^n)
\end{align*}
% [原文 p.115]
成立。这里，从第 1 行到第 2 行的变形中用了 $p_0^{n+1}\circ\pi^{n+1}\circ d_0=0$。

这个 $p=\{p^n\}_{n\in\Z}$ 是拟同构映射，也就是说，要证明 $H^n(p)\colon H^n(P)\to H^n(X)$ 是全双射。

\para{［$H^n(p)$ 的全射性］}
首先，关注态射 $d_X^n\circ p_0^n=p_0^{n+1}\circ\phi_0^n\colon P_0^n\to X^{n+1}$ 的核，证明 $Z^n(P)\simeq\Ker(d_X^n\circ p_0^n)$。对每个 $n\in\Z$，取 $\iota_{Z(P)}\colon Z^n(P)\to P^n$、$P^n=\bigoplus_{i=0}^\infty P_i^{n+i}$，把其中对应于 $Z^n(P)\to P_i^{n+i}$ 的成分设为
\begin{equation}
  x_i:=\iota_{Z(P)}|_{Z^n(P)\to P_i^{n+i}},\qquad i=0,1,\ldots \tag{5.19}
\end{equation}
把 $x_i\colon Z^n(P)\to P_i^{n+i}$ 与单射 $P_i^{n+i}\to P^n$ 的合成也用同一记号 $x_i\colon Z^n(P)\to P^n$ 表示，则 $\iota_{Z(P)}=x_0+x_1+\cdots$。条件式
\[
  0=d_P^n\circ\iota_{Z(P)}=(d_0+d_1+\cdots)\circ(x_0+x_1+\cdots)\colon Z^n(P)\to P^{n+1}
\]
中，关注对应于 $Z^n(P)\to P_0^{n+1}$ 的成分，可得 $d_1\circ x_0+d_0\circ x_1=0$。把它与 $p_0^{n+1}$ 合成，得
\begin{align*}
  0 &= -(p_0^{n+1}\circ d_0)\circ x_1 = -p_0^{n+1}\circ(d_0\circ x_1)\\
    &= p_0^{n+1}\circ(d_1\circ x_0) = (p_0^{n+1}\circ d_1)\circ x_0\\
    &= (d_X^n\circ p_0^n)\circ x_0.
\end{align*}
由此，$x_0\colon Z^n(P)\to P_0^n$ 经过 $\Ker(d_X^n\circ p_0^n)$ 分解。也就是说，$x_0=\pi_P\circ\iota_{Z(P)}$ 经过 $\Ker(d_X^n\circ p_0^n)$ 分解。

% [Tutor Note: 底本 p.115 该计算第 2 项印作「$p_0^{n+1}\circ(d_0\circ x_1)$」；据作者正誤表（2014-05-20 版）p.115 line 15 订正为「$-p_0^{n+1}\circ(d_0\circ x_1)$」，此处按订正后排印。核对：由 $d_1\circ x_0+d_0\circ x_1=0$ 得 $d_0\circ x_1=-d_1\circ x_0$，故 $p_0^{n+1}\circ(d_0\circ x_1)=-p_0^{n+1}\circ(d_1\circ x_0)$，下一行的等号才成立。又同页 line 19 底本印作「$x_0=\pi_P\circ\iota_P$」，据正誤表订正为「$x_0=\pi_P\circ\iota_{Z(P)}$」，译文已照订正。]

反之，对单射 $y_0\colon\Ker(d_X^n\circ p_0^n)\to P_0^n$，由图表
\[\begin{tikzcd}[row sep=2.8em, column sep=3.4em]
\Ker(d_X^n\circ p_0^n) \arrow[dr,"y_0"'] \\
 & P_0^n \arrow[r,"d_0"] \arrow[d,"p_0^n"] & P_0^{n+1} \arrow[d,"p_0^{n+1}"]\\
 & X^n \arrow[r,"d_X^n"] & X^{n+1}
\end{tikzcd}\]
的交换性得
\[
  0=(d_X^n\circ p_0^n)\circ y_0=(p_0^{n+1}\circ d_1)\circ y_0=p_0^{n+1}\circ(d_1\circ y_0),
\]
并且由 $P_1^n\xrightarrow{d_0}P_0^n\xrightarrow{p_0^n}X^n$ 的正合性，$d_1\circ y_0\colon\Ker(d_X^n\circ p_0^n)\to P_0^{n+1}$ 经过 $P_1^{n+1}$ 分解。把对应的态射记作 $y_1\colon\Ker(d_X^n\circ p_0^n)\to P_1^{n+1}$，并把 $y_1$ 与单射 $P_1^{n+1}\to P^n$ 的合成也用同一记号 $y_1$ 表示，则关于 $(y_0+y_1)\colon\Ker(d_X^n\circ p_0^n)\to P^n$，有
% [原文 p.116]
$d_P\circ(y_0+y_1)\colon\Ker(d_X^n\circ p_0^n)\to P^{n+1}$
在 $P_0^{n+1}$ 上是平凡的。同样地由归纳法可证明存在 $y_i\colon\Ker(d_X^n\circ p_0^n)\to P_i^{n+i}$ 使 $d_P\circ(\sum_{i=0}^\infty y_i)=0$。设在 $i=k$ 处，
\[
  d_P\circ(y_0+y_1+\cdots+y_k)\colon\Ker(d_X^n\circ p_0^n)\to P^{n+1}
\]
在 $P_0^{n+1}\oplus\cdots\oplus P_{k-1}^{n+k}$ 上是平凡的。也就是说，把 $d_P\circ(y_0+y_1+\cdots+y_k)$ 的 $\Ker(d_X^n\circ p_0^n)\to P_i^{n+(i+1)}$ 对应的成分记作 $t_i$，则 $t_0=\cdots=t_{k-1}=0$，且
\[
  d_P\circ(y_0+y_1+\cdots+y_k)=t_k+t_{k+1}+\cdots
\]
成立。由
\[
  0=(d_P)^2\circ(y_0+y_1+\cdots+y_k)=d_P\circ(d_P\circ(y_0+y_1+\cdots+y_k))
\]
出发，关注其中对应于 $\Ker(d_X^n\circ p_0^n)\to P_{k-1}^{n+(k+1)}$ 的成分，得
\[
  d_0\circ t_k=0.
\]
由 $P_{k+1}^{n+(k+1)}\xrightarrow{d_0}P_k^{n+(k+1)}\xrightarrow{d_0}P_{k-1}^{n+(k+1)}$ 的完全性，$t_k$ 经过 $P_k^{n+(k+1)}$ 分解；把对应的态射记作 $y_{k+1}\colon\Ker(d_X^n\circ p_0^n)\to P_k^{n+(k+1)}$，则
\[
  d_P\circ((y_0+y_1+\cdots+y_k)+y_{k+1})\colon\Ker(d_X^n\circ p_0^n)\to P^{n+1}
\]
在 $(P_0^{n+1}\oplus\cdots\oplus P_{k-1}^{n+k})\oplus P_k^{n+(k+1)}$ 上是平凡的。由以上归纳法，证明了存在 $y_i\colon\Ker(d_X^n\circ p_0^n)\to P_i^{n+i}$ 使 $d_P\circ(\sum_{i=0}^\infty y_i)=0$。这意味着 $y_0+y_1+\cdots\colon\Ker(d_X^n\circ p_0^n)\to P^n$ 经过 $Z^n(P)$ 分解。
\[\begin{tikzcd}[row sep=2.8em, column sep=4.2em]
Z^n(P) \arrow[r,"\iota_{Z(P)}"] \arrow[d,dashed] & P^n \arrow[d,"\pi_P"]\\
\Ker(d_X^n\circ p_0^n) \arrow[r,"y_0"] \arrow[ur,"\sum_{i=0}^{\infty}y_i"'] & P_0^n
\end{tikzcd}\]
以上得到同构 $Z^n(P)\simeq\Ker(d_X^n\circ p_0^n)$，特别地，上面构造中的 $x_i$ 与 $y_i$ 可以视为同一。
现在想证明 $H^n(p)$ 的全射性，而实际上可知 $Z^n(p)$ 是全射的。事实上，由阿贝尔范畴具有足够多的投射对象，存在投射对象 $P_Z\in\cP$ 与 $\cC$ 中的满射 $p_Z\colon P_Z\to Z^n(X)$。这里，考虑图表
% [原文 p.117]
\[\begin{tikzcd}[row sep=2.8em, column sep=2.6em]
 & Z^n(P) \arrow[r,"x_0"] \arrow[d,"Z^n(p)"] & P_0^n \arrow[r,"d_1"] \arrow[d,"p_0^n"] & P_0^{n+1} \arrow[d,"p_0^{n+1}"]\\
P_Z \arrow[r,"p_Z"] & Z^n(X) \arrow[r,"\iota_{Z(X)}"] & X^n \arrow[r,"d_X^n"] & X^{n+1}
\end{tikzcd}\]
时，由 $P_Z$ 的投射性与 $p_0^n$ 是满射，态射 $\iota_{Z(X)}\circ p_Z$ 经过 $P_0^n$ 分解；把对应的态射记作 $l'\colon P_Z\to P_0^n$，则
\[
  (d_X^n\circ p^n)\circ l'=0
\]
且由 $Z^n(P)\simeq\Ker(d_X^n\circ p_0^n)$，$l'$ 经过 $Z^n(P)$ 分解。把对应的态射记作 $l\colon P_Z\to Z^n(P)$，则得到以下图表
\[\begin{tikzcd}[row sep=2.8em, column sep=2.6em]
 & Z^n(P) \arrow[r,"x_0"] \arrow[d,"Z^n(p)"] & P_0^n \arrow[r,"d_1"] \arrow[d,"p_0^n"] & P_0^{n+1} \arrow[d,"p_0^{n+1}"]\\
P_Z \arrow[r,"p_Z"] \arrow[ur,dashed,"l"] \arrow[urr,dashed,"l'"] & Z^n(X) \arrow[r,"\iota_{Z(X)}"] & X^n \arrow[r,"d_X^n"] & X^{n+1}
\end{tikzcd}\]
且 $p_Z=Z^n(p)\circ l$。这里因 $p_Z$ 是满射，可知 $Z^n(p)$ 是满射。此时，图表
\[\begin{tikzcd}[row sep=2.4em, column sep=2.6em]
0 \arrow[r] & B^n(P) \arrow[r,"\iota_P"] \arrow[d,"B^n(p)"] & Z^n(P) \arrow[r,"\pi_P"] \arrow[d,"Z^n(p)"] & H^n(P) \arrow[r] \arrow[d,"H^n(p)"] & 0\\
0 \arrow[r] & B^n(X) \arrow[r,"\iota_X"] & Z^n(X) \arrow[r,"\pi_X"] & H^n(X) \arrow[r] & 0
\end{tikzcd}\]
中，由 $\pi_X\circ Z^n(p)=H^n(p)\circ\pi_P$ 是全射，可知 $H^n(p)$ 是全射。\hfill$\square$

% [Tutor Note: 底本 p.117 第三个图式的下排印作 $B^i(X), Z^i(X), H^i(X)$；据作者正誤表（2014-05-20 版）p.117「３つ目の図式」条订正为 $B^n(X), Z^n(X), H^n(X)$（同一图式在 p.119 再次出现，也按此订正），译文已按订正排印。又 (5.20) 直前底本印作 $\iota_{B(X)}:B^n(X)\to Z^n(X)$，据正誤表 p.117「(5.20) の直前」条订正为 $\iota_X:B^n(X)\to Z^n(X)$，译文亦按订正。可核对：$\iota_X$ 即上式图底行 $B^n(X)\hookrightarrow Z^n(X)$ 的单射。又底本 p.117 印作 $(d_X^n\circ p^n)\circ l'=0$，按类型 $l'\colon P_Z\to P_0^n$ 应为 $(d_X^n\circ p_0^n)\circ l'=0$；正誤表未收，译文照底本录入，在此注明。]

\para{［$H^n(p)$ 的单射性］}
首先，取投射对象 $P_W$ 与态射 $p_W\colon P_W\to Z^n(P)$，要证明：当 $Z^n(p)\circ p_W\colon P_W\to Z^n(X)$ 经过 $\iota_X\colon B^n(X)\to Z^n(X)$ 分解时，$p_W$ 必定经过 $\iota_P\colon B^n(P)\to Z^n(P)$ 分解。也就是说，证明图表
\begin{equation}
\begin{tikzcd}[row sep=3.4em, column sep=2.8em]
 & & P_W \arrow[d,"p_W"] \arrow[dll,"s"'] \arrow[ddl,"q"] & \\
P^{n-1} \arrow[r,"\pi_{B(P)}"] \arrow[d,"p_0^{n-1}"] & B^n(P) \arrow[r,"\iota_P"] \arrow[d,"B^n(p)"] & Z^n(P) \arrow[r,"\iota_{Z(P)}"] \arrow[d,"Z^n(p)"] & P^n \arrow[d,"p^n"]\\
X^{n-1} \arrow[r,"\pi_{B(X)}"] & B^n(X) \arrow[r,"\iota_X"] & Z^n(X) \arrow[r,"\iota_{Z(X)}"] & X^n
\end{tikzcd} \tag{5.20}
\end{equation}
中，当 $Z^n(p)\circ p_W=\iota_X\circ q$ 时，必定有 $\iota_{Z(P)}\circ p_W=d_P\circ s$ 成立。此时，因 $\iota_{Z(P)}$ 是单射，得 $p_W=\iota_P\circ\pi_{B(P)}\circ s$。

% [Tutor Note: 底本 p.117 此后对 $P_W$ 连了三条箭：竖直的 $p_W$、指向 $P^{n-1}$ 的 $s$、指向 $B^n(X)$ 的 $q$。核对：$q\colon P_W\to B^n(X)$ 来自 $Z^n(p)\circ p_W$ 经 $\iota_X$ 的唯一分解（$\iota_X$ 单射）；$s\colon P_W\to P^{n-1}$ 是用 $P_W$ 的投射性把 $\iota_{Z(P)}\circ p_W$ 沿满射 $P^{n-1}\twoheadrightarrow B^n(P)$ 提升所得。又底本左端节点印作 $P^{n-1}$（漏下标 0），译文照印。]

% [原文 p.118]
在图式 (5.20) 中，射
\[
  \iota_{Z(X)}\circ Z^n(p)\circ p_W=p^n\circ\iota_{Z(P)}\circ p_W\colon P_W\to X^n
\]
若经过 $B(X)$ 分解，则经过 $X^{n-1}$ 分解。进一步，由于 $p_0^{n-1}$ 是满射，由 $P_W$ 的投射性，存在态射 $s_0\colon P_W\to P_0^{n-1}$，使得
\[
  p_0^n\circ x_0\circ p_W=p^n\circ\iota_{Z(P)}\circ p_W=(d_X^{n-1}\circ p_0^{n-1})\circ s_0=(p_0^n\circ d_1)\circ s_0.
\]
\[\begin{tikzcd}[row sep=3.0em, column sep=3.2em]
P_W \arrow[r,"p_W"] \arrow[dr,dashed,"s_0"'] \arrow[ddr] & Z^n(P) \arrow[r,"\iota_{Z(P)}"] \arrow[dr,"x_0"'] & P^n \arrow[d,"\pi"]\\
 & P_0^{n-1} \arrow[r,"d_1"] \arrow[d,"p_0^{n-1}"] & P_0^n \arrow[d,"p_0^n"]\\
 & X^{n-1} \arrow[r,"d_X^{n-1}"] & X^n
\end{tikzcd}\]
由此，$p_0^n\circ(x_0\circ p_W-d_1\circ s_0)=0$，故由 $P_1^n\xrightarrow{d_0}P_0^n\xrightarrow{p_0^n}X^n$ 的完全性，存在 $s_1\colon P_W\to P_1^n$ 使 $x_0\circ p_W-d_1\circ s_0=d_0\circ s_1$。也就是说，对于 $s_0+s_1\colon P_W\to P^{n-1}$，有
\[
  x_0\circ p_W=d_1\circ s_0+d_0\circ s_1=(d_P\circ(s_0+s_1))|_{P_W\to P_0^n}
\]
成立。由于 $\iota_{Z(P)}=x_0+x_1+\cdots$，也可以写作 $(\iota_{Z(P)}\circ p_W-d_P(s_0+s_1))|_{P_W\to P_0^n}=0$。

以下用归纳法，对 $i=0,1,\ldots$ 证明：存在满足
\begin{equation}
\begin{gathered}
  (S_i)|_{P_W\to(P_0^n\oplus\cdots\oplus P_i^{n+i})}=0,\\
  S_i:=\iota_{Z(P)}\circ p_W-d_P\circ(s_0+s_1+\cdots+s_{i+1})
\end{gathered} \tag{5.21}
\end{equation}
的 $s_i\colon P_W\to P_i^{n+i-1}$。设 $i=k$ 时 (5.21) 成立。也就是说，$S_k$ 在 $P_0^n,\ldots,P_k^{n+k}$ 上为零，但在 $P_{k+1}^{n+k+1},\ldots$ 上未必平凡。此时，$d_P\circ S_k$ 在 $P_0^{n+1},\ldots,P_{k-1}^{n+k}$ 上显然是平凡的；而由 $S_k$ 的构造本身有 $0=d_P\circ S_k$，故把其在 $P_k^{n+k+1}$ 上的条件写出，便得
\[
  0=(d_P\circ S_k)|_{P_W\to P_k^{n+k+1}}=d_0\circ(S_k)|_{P_W\to P_{k+1}^{n+k+1}}.
\]

% [Tutor Note: 底本 p.118 此处印作「$P_1^n\xrightarrow{d_0}P_0^n\xrightarrow{d_0}X^n$ の完全性より」；第二个箭头按对象应为 $p_0^n\colon P_0^n\to X^n$（与 p.115 同一论证处的印法一致），作者正誤表未收此条，译文照底本印 $d_0$ 并在此注明。可核对：所用正合性即 $0\leftarrow X^n\xleftarrow{p_0^n}P_0^n\xleftarrow{d_0}P_1^n$ 的正合性。]

% [原文 p.119]
\[\begin{tikzcd}[row sep=2.6em, column sep=2.2em]
P_W \arrow[drrrr,"(S_k)|_{P_W\to P_{k+1}^{n+k+1}}"] \arrow[ddrrr,"0"] \arrow[dddrr,"0"] \arrow[ddddr,"0"] & & & & \\
& & & & P_{k+1}^{n+k+1} \arrow[d,"d_0"]\\
& & & P_k^{n+k} \arrow[r,"d_1"] \arrow[d,"d_0"] & P_k^{n+k+1}\\
& & P_1^{n+1} \arrow[r,"d_1"] \arrow[d,"d_0"] & \cdots & \\
& P_0^n \arrow[r,"d_1"] & P_0^{n+1} & &
\end{tikzcd}\]
因此，$(S_k)|_{P_W\to P_{k+1}^{n+k+1}}$ 经过 $P_{k+2}^{n+k+1}$ 分解；把对应的态射记作 $s_{k+2}\colon P_W\to P_{k+2}^{n+k+1}$，并令
\[
  S_{k+2}:=\iota_{Z(P)}\circ p_W-d_P\circ((s_0+s_1+\cdots+s_{k+1})+s_{k+2}),
\]
便得到 $(S_{k+2})|_{P_W\to((P_0^n\oplus\cdots\oplus P_k^{n+k})\oplus P_{k+1}^{n+k+1})}=0$。于是归纳法完成；令
\[
  s:=s_0+s_1+\cdots,
\]
则 $\iota_{Z(P)}\circ p_W=d_P\circ s$，也就是说，可知 $\iota_{Z(P)}\circ p_W$ 经过 $P^{n-1}$ 分解。这意味着 $\iota_{Z(P)}\circ p_W$ 经过 $B^n(P)$ 分解。因此，$p_W$ 经过 $B^n(P)$ 分解。

% [Tutor Note: 底本 p.119 上图右端节点印作 $P_0^{n+k+1}$；据作者正誤表（2014-05-20 版）p.119「１つ目の図式」条订正为 $P_k^{n+k+1}$，此处按订正后重画。核对：由 $d_1$ 的双次数 $(\Delta n,\Delta i)=(1,0)$ 知 $d_1\colon P_k^{n+k}\to P_k^{n+k+1}$；右侧竖直箭头 $d_0$ 则把 $P_{k+1}^{n+k+1}$ 送到 $P_k^{n+k+1}$（$d_0$ 的 $(\Delta n,\Delta i)=(0,-1)$）。图中 $P_W$ 到各层的三条「$0$」箭正是归纳假设所给 $S_k=0$ 的成分。]

现在可以证明 $H^n(p)\colon H^n(P)\to H^n(X)$ 的单射性。也就是说，要证明：$\cC$ 中的态射 $w\colon W\to H^n(P)$ 若满足 $H^n(p)\circ w=0$，则恒有 $w=0$。由阿贝尔范畴 $\cC$ 具有足够多的投射对象，存在投射对象 $P_W\in\cP$ 与满射 $p_W\colon P_W\to W$。
\[\begin{tikzcd}[row sep=2.6em, column sep=2.4em]
 & & P_W \arrow[ddl,dashed,"w''"'] \arrow[dddl,dashed,"w'"] \arrow[d,"p_W"] & & \\
 & & W \arrow[d,"w"] & & \\
0 \arrow[r] & B^n(P) \arrow[r,"\iota_P"] \arrow[d,"B^n(p)"] & Z^n(P) \arrow[r,"\pi_P"] \arrow[d,"Z^n(p)"] & H^n(P) \arrow[r] \arrow[d,"H^n(p)"] & 0\\
0 \arrow[r] & B^n(X) \arrow[r,"\iota_X"] & Z^n(X) \arrow[r,"\pi_X"] & H^n(X) \arrow[r] & 0
\end{tikzcd}\]
$P_W$ 是投射的，且 $\pi_P\colon Z^n(P)\to H^n(P)$ 是全射，故 $w\circ p_W$
% [原文 p.120]
经过 $\pi_P\colon Z^n(P)\to H^n(P)$ 分解，得到满足
\[
  w\circ p_W=\pi_P\circ w'
\]
的态射 $w'\colon P_W\to Z^n(P)$。关于从 $P_W$ 到 $H^n(X)$ 的合成射，有
\begin{align*}
  0 &= H^n(p)\circ(w\circ p_W)\\
    &= H^n(p)\circ(\pi_P\circ w')=(H^n(p)\circ\pi_P)\circ w'\\
    &= (\pi_X\circ Z^n(p))\circ w'=\pi_X\circ(Z^n(p)\circ w')
\end{align*}
由此可知，$Z^n(p)\circ w'$ 经过 $\iota_X\colon B^n(X)\to Z^n(X)$ 分解。

这里，由上面的论证，$w'$ 经过 $\iota_P\colon B^n(P)\to Z^n(P)$ 分解；把对应的态射记作 $w''\colon P_W\to B^n(P)$，则 $w'=\iota_P\circ w''$。由以上，
\[
  w\circ p_W=\pi_P\circ\iota_P\circ w''=0,
\]
而 $p_W$ 是全射，故 $w=0$，这就证明了 $H^n(p)$ 的单射性。\hfill$\square$

% [Tutor Note: 底本 p.120 line 10 印作「$w'=\iota_{B(P)}\circ w''$」；据作者正誤表（2014-05-20 版）p.120 条订正为「$w'=\iota_P\circ w''$」，此处按订正后排印。核对：$\iota_P$ 是 $B^n(P)\to Z^n(P)$ 的单射（p.117 图式第三图的左段），而 $w''$ 的陪域是 $B^n(P)$。又 p.119/p.120 的图示中 $w'$ 与 $w''$ 皆为虚线箭，$w''$ 指向 $B^n(P)$、$w'$ 指向 $Z^n(P)$，二者关系为 $w'=\iota_P\circ w''$。]

\begin{remark}{5.41}
定理 5.36 是场论中 Batalin–Vilkovisky 形式的基本定理，对应于把称为同调扰动理论（homological perturbation theory）的结论 [19, p178, Theorem 8.3] 在投射对象的双系列情形下改写所得的结果。
\end{remark}

\subsection{关于投射分解的另一个基本定理}

在补題 5.34 中，对 $\cC$ 的态射 $f\colon X\to Y$，把 $X\in\cC, Y\in\cC$ 视为复形、把 $f\colon X\to Y$ 视为复形的态射，并把对应于投射分解
\[
  0\leftarrow X\xleftarrow{p_0}P_0\xleftarrow{p_1}P_1\leftarrow\cdots,\qquad
  0\leftarrow Y\xleftarrow{q_0}Q_0\xleftarrow{q_1}Q_1\leftarrow\cdots
\]
的由投射对象构成的复形记作 $P, Q$（cf. (4.18)），则补題 5.34 的断言等价于：存在使以下图表交换的复形态射 $\phi\colon P\to Q$，且除同伦外唯一。
\[\begin{tikzcd}[row sep=2.8em, column sep=4.0em]
P \arrow[r,"\phi"] \arrow[d,"p"] & Q \arrow[d,"q"]\\
X \arrow[r,"f"] & Y
\end{tikzcd}
\qquad
p:=p_0,\qquad q:=q_0,\qquad
\phi^i=\begin{cases}\phi_{-i}, & i\le 0,\\ 0, & i>0.\end{cases}\]
考虑它的推广。

\begin{theorem}{5.42}
对两个上方有界复形 $X, Y\in\Comp^-(\cC)$ 与复形态射 $f\colon X\to Y$，设给定了由投射对象构成的上方有界复形 $P, Q\in\Comp^-(\cP)$ 与复形的拟同构映射 $p\colon P\to X$、$q\colon Q\to Y$。此时，
\begin{enumerate}[label=(\roman*)]
  \item 存在复形态射 $\tilde f\colon P\to Q$，使以下图表
\[\begin{tikzcd}[row sep=2.8em, column sep=4.0em]
P \arrow[r,"\tilde f"] \arrow[d,"p"] & Q \arrow[d,"q"]\\
X \arrow[r,"f"] & Y
\end{tikzcd}\]
 在 $K^-(\cC)$ 中可换。
% [原文 p.121]
  \item 这样的射 $\tilde f\colon P\to Q$ 在同伦范畴 $K^-(\cP)$ 中唯一。
\end{enumerate}
\end{theorem}

% [Tutor Note: 定理 5.42 的证明一直延续到 p124 的 □ 才结束，跨 chunk 边界；与命題 5.38 同样按原文体例用 \para 手工起头，末端的 □ 在 p124 处补上（排版效果与 proof 环境一致）。]

\para{定理 5.42 的证明。}
(i) 对每个 $i$，归纳地构造满足
\begin{equation}
  d_Q^i\circ\tilde f^i=\tilde f^{i+1}\circ d_P^i, \tag{5.22}
\end{equation}
\begin{equation}
  q^i\circ\tilde f^i-f^i\circ p^i=d_Y^{i-1}\circ h^i+h^{i+1}\circ d_P^i \tag{5.23}
\end{equation}
的 $\tilde f^i\colon P^i\to Q^i$、$h^i\colon P^i\to Y^{i-1}$。
\[\begin{tikzcd}[row sep=2.6em, column sep=1.5em]
\cdots \arrow[rr] & & Q^i \arrow[rr,"d_Q^i"] & & Q^{i+1} \arrow[rr,"d_Q^{i+1}"] & & Q^{i+2} \arrow[r] & \cdots \\
\cdots \arrow[r] & P^i \arrow[rr,"d_P^i"] \arrow[ur,"\tilde f^i"] \arrow[dd,"p^i"] & & P^{i+1} \arrow[rr,"d_P^{i+1}"] \arrow[ur,"\tilde f^{i+1}"] \arrow[dd,"p^{i+1}"] & & P^{i+2} \arrow[ur,"\tilde f^{i+2}"] \arrow[dd,"p^{i+2}"] \arrow[r] & {} & \\
& & Y^i \arrow[rr,"d_Y^i"] & & Y^{i+1} \arrow[rr,"d_Y^{i+1}"] & & Y^{i+2} \arrow[r] & \cdots \\
\cdots \arrow[r] & X^i \arrow[rr,"d_X^i"] \arrow[ur,"f^i"] & & X^{i+1} \arrow[rr,"d_X^{i+1}"] \arrow[ur,"f^{i+1}"] & & X^{i+2} \arrow[ur,"f^{i+2}"] \arrow[r] & {} &
\end{tikzcd}\]
首先，对足够大的 $i>N$，设 $P^i=Q^i=X^i=Y^i=0$，则 $p^i, q^i, f^i$ 全为零射，取 $\tilde f^i=0, h^i=0$ 即满足条件 (5.23)；且由 $d_P^i=0, d_Q^i=0$ 也满足 (5.22)。
其次，假设对 $i\ge k+1$ 存在满足 (5.22),(5.23) 的 $\tilde f^i, h^i$，构造 $i=k$ 时满足 (5.22),(5.23) 的 $\tilde f^k, h^k$。
首先构造满足 (5.22) 的 $\tilde f^k$。关注图表
\[\begin{tikzcd}[row sep=2.6em, column sep=3.0em]
Q^k \arrow[r,"d_Q^k"] & Q^{k+1} \arrow[r,"d_Q^{k+1}"] & Q^{k+2}\\
P^k \arrow[r,"d_P^k"] & P^{k+1} \arrow[u,"\tilde f^{k+1}"] \arrow[r,"d_P^{k+1}"] & P^{k+2} \arrow[u,"\tilde f^{k+2}"]
\end{tikzcd}\]
时，由
\begin{align*}
  d_Q^{k+1}\circ(\tilde f^{k+1}\circ d_P^k)
  &= (d_Q^{k+1}\circ\tilde f^{k+1})\circ d_P^k = (\tilde f^{k+2}\circ d_P^{k+1})\circ d_P^k\\
  &= \tilde f^{k+2}\circ(d_P^{k+1}\circ d_P^k)=0
\end{align*}
可知，射 $\tilde f^{k+1}\circ d_P^k\colon P^k\to Q^{k+1}$ 经过 $Z^{k+1}(Q)$ 分解。把对应的态射记作 $u\colon P^k\to Z^{k+1}(Q)$，考虑以下图表。
% [原文 p.122]
\[\begin{tikzcd}[row sep=2.6em, column sep=2.6em]
 & & P^k \arrow[d,"u"] & & \\
0 \arrow[r] & B^{k+1}(Q) \arrow[r,"\iota_Q"] \arrow[d,"B^{k+1}(q)"] & Z^{k+1}(Q) \arrow[r,"\pi_Q"] \arrow[d,"Z^{k+1}(q)"] & H^{k+1}(Q) \arrow[r] \arrow[d,"H^{k+1}(q)"] & 0\\
0 \arrow[r] & B^{k+1}(Y) \arrow[r,"\iota_Y"] & Z^{k+1}(Y) \arrow[r,"\pi_Y"] & H^{k+1}(Y) \arrow[r] & 0
\end{tikzcd}\]
在该图表中，合成 $Z^{k+1}(q)\circ u$ 经过 $B^{k+1}(Y)$ 分解。实际上，把 $i=k+1$ 时的 (5.23) 的两边与 $d_P^k$ 合成，得
\[
  (q^{k+1}\circ\tilde f^{k+1}-f^{k+1}\circ p^{k+1})\circ d_P^k=d_Y^k\circ h^{k+1}\circ d_P^k,
\]
而由于 $f, p$ 是复形的态射，有 $f^{k+1}\circ p^{k+1}\circ d_P^k=d_Y^k\circ f^k\circ p^k$，故
\[
  q^{k+1}\circ\tilde f^{k+1}\circ d_P^k=d_Y^k\circ(f^k\circ p^k+h^{k+1}\circ d_P^k).
\]
这意味着 $Z^{k+1}(q)\circ u$ 经过 $B^{k+1}(Y)$ 分解。
由此，$\pi_Y\circ Z^{k+1}(q)\circ u=0$；又因 $H^{k+1}(q)$ 是同构态射，故 $\pi_Q\circ u=0$。因此，可知 $u$ 经过 $B^{k+1}(Q)$ 分解。由 $P$ 的投射性与满射 $Q^k\to B^{k+1}(Q)$ 的存在，$u$ 经过 $Q^k$ 分解；也就是说，可知在 $i=k$ 处存在满足 (5.22) 的 $\tilde f^k$。

% [Tutor Note: 底本 p.122 line 11 印作「全射 $Q^{k+1}\to B^{k+1}(Q)$ の存在により，$u$ が $Q^{k+1}$ を経由し」；据作者正誤表（2014-05-20 版）p.122 条订正为「全射 $Q^k\to B^{k+1}(Q)$ の存在により，$u$ が $Q^k$ を経由し」，此处按订正后排印。核对：$B^{k+1}(Q)=\Img(d_Q^k)$，故满射是 $Q^k\twoheadrightarrow B^{k+1}(Q)$（由 $d_Q^k\colon Q^k\to Q^{k+1}$ 诱导），而不是从 $Q^{k+1}$ 出发。]

取这样的一个 $\tilde f^k$，记为 $\tilde f_0^k$。此时由
\[
  d_Y^k\circ(q^k\circ\tilde f_0^k-f^k\circ p^k-h^{k+1}\circ d_P^k)=0,
\]
令
\[
  \varphi^k:=q^k\circ\tilde f_0^k-f^k\circ p^k-h^{k+1}\circ d_P^k,
\]
则 $\varphi^k\colon P^k\to Y^k$ 经过 $Z^k(Y)$ 分解。要证明 $\varphi^k$ 进一步经过 $Y^{k-1}\oplus Z^k(Q)$ 分解。为此，只需证明由 $q^k\colon Q^k\to Y^k$ 诱导的态射 $Z^k(Q)\to Z^k(Y)$ 与 $B^k(Y)\to Z^k(Y)$ 所构成的态射 $Z^k(Q)\oplus B^k(Y)\to Z^k(Y)$ 是全射。实际上，对于某个态射 $w\colon Z^k(Y)\to W$，考虑以下图表
\[\begin{tikzcd}[row sep=2.4em, column sep=2.6em]
0 \arrow[r] & B^k(Q) \arrow[r,"\iota_Q"] \arrow[d,"B^k(q)"] & Z^k(Q) \arrow[r,"\pi_Q"] \arrow[d,"Z^k(q)"] & H^k(Q) \arrow[r] \arrow[d,"H^k(q)"] & 0\\
0 \arrow[r] & B^k(Y) \arrow[r,"\iota_Y"] & Z^k(Y) \arrow[r,"\pi_Y"] \arrow[d,"w"] & H^k(Y) \arrow[r] & 0\\
 & & W & &
\end{tikzcd}\]
若 $w\circ(Z^k(q)\oplus\iota_Y)=0$，则
% [原文 p.123]
\begin{equation}
  w\circ Z^k(q)=0, \tag{5.24}
\end{equation}
\begin{equation}
  w\circ\iota_Y=0 \tag{5.25}
\end{equation}
成立。由 (5.25)，$w$ 经过 $H^k(Y)$ 分解。记作 $\tilde w\colon H^k(Y)\to W$。其次，由 (5.24)，
\[
  0=w\circ Z^k(q)=\tilde w\circ H^k(q)\circ\pi_Q,
\]
而 $H^k(q)\circ\pi_Q$ 是全射，故这意味着 $\tilde w=0$。由 $w=\tilde w\circ\pi_Y=0\circ\pi_Y=0$，证明了 $Z^k(q)\oplus\iota_Y$ 是全射。

% [Tutor Note: 底本 p.122 下方图式与 p.123 line 1,5,7,9 把 $B^k(q),Z^k(q),H^k(q)$ 印作 $B^k(\varphi),Z^k(\varphi),H^k(\varphi)$；据作者正誤表（2014-05-20 版）p.122「下の図式」条与 p.123 条，一律订正为 $q$ 的版本（$B^k(q),Z^k(q),H^k(q)$），此处按订正后排印。核对：该图式是 $q\colon Q\to Y$ 的复形态射所诱导的 $B/Z/H$ 长正合列之间的比较图，与 $\varphi^k$（它只是 $P^k\to Y^k$ 的态射）无关。]

现在，由 $Z^k(q)\oplus\iota_Y$ 的全射性，$\varphi^k\colon P^k\to Y^k$ 经过 $Y^{k-1}\oplus Z^k(Q)$ 分解。又因存在单射 $Z^k(Q)\to Q^k$，故 $\varphi^k$ 经过 $Y^{k-1}\oplus Q^k$ 分解。把对应的态射表示为
\[
  h^k\oplus g^k\colon P^k\to Y^{k-1}\oplus Q^k,
\]
则可知这个 $h^k$ 与 $\tilde f^k:=\tilde f_0^k+g^k$ 满足 (5.23)。把 $\tilde f_0^k$ 改为 $\tilde f^k$ 后条件 (5.22) 仍保持，这一点由 $g^k$ 经过 $Z^k(Q)$ 分解而得到保证。于是归纳法完成。\hfill$\square$

(ii) 设给出了 (i) 的两个复形态射，记为 $\tilde f,\tilde g\colon P\to Q$。对足够大的 $i>N$，设 $P^i=Q^i=X^i=Y^i=0$，则 $f^i$ 为零射，且唯一。
其次，设在 $i\ge k$ 处，存在满足
\begin{equation}
  \tilde f^i-\tilde g^i=d_Q^{i-1}\circ\tilde h^i+\tilde h^{i+1}\circ d_P^i \tag{5.26}
\end{equation}
的 $\tilde h^i\colon P^i\to Q^{i-1}$（$i=k,k+1,k+2,\ldots$）。（特别地，对足够大的 $i$ 有 $\tilde h^i=0$。）此时要证明：在 $i=k-1$ 的情形，存在满足 (5.26) 的 $\tilde h^{k-1}\colon P^{i-1}\to Q^{i-2}$。

% [Tutor Note: 底本 p.123 line −11 在此重复印出「$\tilde h^{k-1}:\tilde h^{k-1}:P^{i-1}\to Q^{i-2}$」；据作者正誤表（2014-05-20 版）p.123 条，重复的冒号已删去，此处照订正后排印（保留底本的上标 $i-1,i-2$；按上下文当为 $k-1,k-2$ 之误，正誤表未及，译文照印）。]

\begin{align*}
  d_Q^{k-1}\circ(\tilde f^{k-1}-\tilde g^{k-1})
  &= (\tilde f^k-\tilde g^k)\circ d_P^{k-1}\\
  &= (d_Q^{k-1}\circ\tilde h^k+\tilde h^{k+1}\circ d_P^k)\circ d_P^{k-1}
   = d_Q^{k-1}\circ\tilde h^k\circ d_P^{k-1}
\end{align*}
由此，$d_Q^{k-1}\circ(\tilde f^{k-1}-\tilde g^{k-1}-\tilde h^k\circ d_P^{k-1})=0$ 成立。现在令
\[
  \psi^{k-1}:=\tilde f^{k-1}-\tilde g^{k-1}-\tilde h^k\circ d_P^{k-1}.
\]
由于 $d_Q^{k-1}\circ\psi^{k-1}=0$，$\psi^{k-1}\colon P^{k-1}\to Q^{k-1}$ 经过 $Z^{k-1}(Q)$ 分解。另一方面，
\begin{align*}
  q^{k-1}\circ\psi^{k-1}
  &= q^{k-1}\circ(\tilde f^{k-1}-\tilde g^{k-1}-\tilde h^k\circ d_P^{k-1})\\
  &= q^{k-1}\circ(\tilde f^{k-1}-\tilde g^{k-1})-(q^{k-1}\circ\tilde h^k)\circ d_P^{k-1},
\end{align*}
由此存在 $\bar h^{k-1}$ 使 $q^{k-1}\circ\psi^{k-1}=d_Y^{k-2}\circ\bar h^{k-1}$。也就是说，$q^{k-1}\circ\psi^{k-1}$
% [原文 p.124]
经过 $B^{k-1}(Y)$ 分解。

利用这些，证明 $\psi^{k-1}$ 经过 $B^{k-1}(Q)$ 分解。首先，由 $\psi^{k-1}$ 经过 $Z^{k-1}(Q)$ 分解，把对应的态射记作 $\tilde\psi^{k-1}\colon P^{k-1}\to Z^{k-1}(Q)$，则得到以下图表。
\[\begin{tikzcd}[row sep=2.4em, column sep=2.4em]
P^{k-1} \arrow[drr,"\tilde\psi^{k-1}"] & & & & \\
0 \arrow[r] & B^{k-1}(Q) \arrow[r,"\iota_Q"] \arrow[d,"B^{k-1}(q)"] & Z^{k-1}(Q) \arrow[r,"\pi_Q"] \arrow[d,"Z^{k-1}(q)"] & H^{k-1}(Q) \arrow[r] \arrow[d,"H^{k-1}(q)"] & 0\\
0 \arrow[r] & B^{k-1}(Y) \arrow[r,"\iota_Y"] & Z^{k-1}(Y) \arrow[r,"\pi_Y"] & H^{k-1}(Y) \arrow[r] & 0
\end{tikzcd}\]
这里，横向的行是短正合列。$Z^{k-1}(q)\circ\tilde\psi^{k-1}$ 经过 $B^{k-1}(Y)$ 分解，故 $\pi_Y\circ Z^{k-1}(q)\circ\tilde\psi^{k-1}=0$；又 $H^{k-1}(q)$ 是同构，故得 $\pi_Q\circ\tilde\psi^{k-1}=0$。因此，$\tilde\psi^{k-1}$ 经过 $B^{k-1}(Q)$ 分解。也就是说，$\psi^{k-1}$ 也经过 $B^{k-1}(Q)$ 分解。
$P^{k-1}$ 是投射的，且对由 $d_Q^{k-2}\colon Q^{k-2}\to Q^{k-1}$ 得到的满射 $Q^{k-2}\to B^{k-1}(Q)$，态射 $P^{k-1}\to B^{k-1}(Q)$ 经过 $Q^{k-2}$ 分解。因此，存在满足
\[
  d_Q^{k-2}\circ\tilde h^{k-1}=\psi^{k-1}=\tilde f^{k-1}-\tilde g^{k-1}-\tilde h^k\circ d_P^{k-1}
\]
的 $\tilde h^{k-1}\colon P^{k-1}\to Q^{k-2}$，归纳法完成。\hfill$\square$

\begin{corollary}{5.43}
对复形 $X\in\Comp^-(\cC)$，由投射对象构成且带有拟同构映射 $p\colon P\to X$ 的复形 $P\in\Comp^-(\cP)$，若存在则在同伦范畴 $K^-(\cP)$ 中除同构外唯一确定。
\end{corollary}

\begin{proof}
在补題 5.42 中取 $Y=X$、$f=\id_X$ 的情形。此时存在射 $\tilde f\colon P\to Q$，使 $q\circ\tilde f$ 与 $p$ 在 $\Comp(\cC)$ 中同伦。同样地，存在射 $\tilde f'\colon Q\to P$，使 $p\circ\tilde f'$ 与 $q$ 在 $\Comp(\cC)$ 中同伦。
另一方面，这次在补題 5.42 中取 $Y=X$、$f=\id_X$、$Q=P$、$q=p$ 的情形。此时，作为从复形 $P$ 到 $P$ 的射，有恒等映射 $\id_P$ 与合成 $\tilde f'\circ\tilde f$，它们互相同伦。同样地，合成 $\tilde f\circ\tilde f'\colon Q\to Q$ 与 $\id_Q$ 同伦。以上证明了命题。
\end{proof}

% [Tutor Note: 底本 p.124 系 5.43 的证明中把 5.42 称作「補題 5.42」；但 5.42 在本书中是定理（定理 5.42）。作者正誤表未收此条，译文照底本作「补題 5.42」并在此注明。]

\begin{corollary}{5.44}
对由投射对象构成的复形 $P, Q\in\Comp^-(\cP)$，若存在从 $P$ 到 $Q$ 的拟同构映射，则在同伦范畴 $K^-(\cP)$ 中 $P\simeq Q$。
\end{corollary}

\begin{proof}
现在存在两个拟同构映射 $\id_Q\colon Q\to Q$ 与 $P\to Q$，在系 5.43 中取 $X=Q$ 即可推出 $P\simeq Q$。
\end{proof}

\subsection{结论}

由此可知，投射对象的复形的同伦范畴 $K^-(\cP)$ 已经完成局部化。也就是说，可得以下结论。


